\documentclass[10pt,a4paper]{amsart}
\usepackage[margin=19mm]{geometry}
\usepackage{amsmath,amssymb,mathtools}
\usepackage[scr=rsfs,cal=euler]{mathalfa}
\usepackage{xfrac}
\usepackage[unicode,hidelinks,bookmarks=false]{hyperref}
\newcommand{\ie}{i.e.\ }
\newcommand{\cf}{cf.\ }
\newcommand{\resp}{respectively\ }
\newcommand{\Subsection}{\subsection}
\providecommand{\nobreakdash}{\nobreak\hbox{-}}
\setlength{\emergencystretch}{2em}
\multlinegap0pt
\usepackage{amssymb}
\usepackage{mathtools}
\usepackage{tikz}
\usepackage{algorithm}\usepackage{algpseudocode}
\usepackage[shortlabels]{enumitem}
\usepackage{xcolor}
\newtheorem{proposition}{Proposition}
\newtheorem{theorem}{Theorem}
\title{Vertex connectivity of the nonzero nonunit core of the comaximal graph of $\mathbb Z_n$.}
\author{Bilal Ahmad Rather}
\date{Source: arXiv:2605.01565v1 (CC BY 4.0)}
\begin{document}
\maketitle

\noindent\emph{Source and licence.} Vertex connectivity of the nonzero nonunit core of the comaximal graph of $\mathbb Z_n$.. Version arXiv:2605.01565v1 (CC BY 4.0), distributed by arXiv under the Creative Commons Attribution 4.0 International licence, http://creativecommons.org/licenses/by/4.0/, as recorded in the arXiv licence metadata for this exact version and shown on the abstract page https://arxiv.org/abs/2605.01565v1. Full licence notice: \texttt{licenses/arxiv-2605.01565-CC-BY-4.0.txt}.\par\smallskip

\noindent\textit{Editorial scope.} Counted unit: the article's theorem thm:mindegree (source0.tex lines 363--372). The article's complete original proof of it is delivered with it (source0.tex lines 374--400, one \texttt{proof} environment of 27 lines). No mathematics is written, repaired or re-derived: the statement and the proof are the article's own lines, in the article's own notation, and no retained line is substituted or re-flowed.  Retained uncounted as the source-local context the proof needs: lines 299--305.  Nothing was omitted from any retained span, so the package carries no omission record.  No retained line contains a citation, so the wrapper carries no bibliography and no bibliography entry.  No retained line contains a figure environment, an image-inclusion command or drawing code, so no image is delivered and the package declares no asset dependency.  Editorial formatting only, in addition: an AMS \texttt{amsart} preamble that restates the article's own declarations and macros used by the retained lines, namely the environments \texttt{proposition}, \texttt{theorem} on separate counters, together with the standard packages those lines need.  The retained environments print in this document as Proposition 1, Theorem 1 in order of appearance, whereas the article prints the same items with the numbering of its own full document.  The complete native source of the article as distributed in the arXiv e-print for this exact version is bundled unchanged as \texttt{sources/199644/source0.tex}.
	\begin{proposition}\label{prop:neighborhood}
		If $x\in X_S$, then
		$$
		N(x)=\bigsqcup_{\emptyset\neq T\subseteq S^c} X_T.
		$$
		In particular, the neighborhood depends only on $S$.
	\end{proposition}

	\begin{theorem}\label{thm:mindegree}
		The minimum degree of $G_2$ is
		$$
		\delta(G_2)=\prod_{i=1}^{m-1}(p_i-1),
		$$
		and this value is attained precisely by the vertices in the class
		$$
		X_{[m]\setminus\{m\}}.
		$$
	\end{theorem}

	\begin{proof}
		Let $x\in X_S$, and and let $J=S^c$. Then by Proposition \ref{prop:neighborhood}, we have
		$$
		N(x)=\bigsqcup_{\emptyset\neq T\subseteq J}X_T.
		$$
		Hence, $N(x)$ contains each singleton class $X_{\{j\}}$ with $j\in J$. Therefore, we obtain
		$$
		\deg(x)\geq \sum_{j\in J}|X_{{j}}|.
		$$
		Now, $ |X_{\{j\}}|=\prod_{i\neq j}(p_i-1). $ As $p_1<\cdots<p_m$, the smallest among these singleton-class sizes is
		$$
		|X_{\{m\}}|=\prod_{i=1}^{m-1}(p_i-1).
		$$
		 If $|J|\geq 2$, then $\deg(x)$ contains at least two disjoint singleton layers in its neighborhood, so we have
		$$
		\deg(x)\geq |X_{\{j_1\}}|+|X_{\{j_2\}}|>|X_{\{m\}}|
		$$
		for any distinct $j_1,j_2\in J$. Thus,  no such vertex has minimum degree. If $|J|=1$, say $J=\{j\}$, then $S=[m]\setminus\{j\}$. In this case, the only nonempty subset of $J$ is $\{j\}$ itself, so Proposition \ref{prop:neighborhood} implies that $N(x)=X_{\{j\}},$ 	and therefore, we have
		$$
		\deg(x)=|X_{\{j\}}|=\prod_{i\neq j}(p_i-1).
		$$
		This is minimized exactly when $j=m$, and hence
		$$
		\delta(G_2)=|X_{\{m\}}|=\prod_{i=1}^{m-1}(p_i-1),
		$$
		with equality holds precisely on the class $X_{[m]\setminus\{m\}}$.
	\end{proof}

\end{document}
